Blockchain provenance for additive manufacturing (AM) is only as
trustworthy as the last-mile bridge that converts physical process data
into on-chain commitments---an instance of the oracle problem that
existing AM provenance work leaves largely unaddressed.  We present a
threat-model-driven attestation architecture for the printer-to-ledger
path: a four-tier trust ladder (T0 na\"ive submission, T1 edge-signed
telemetry bundles, T2 on-chain device registry, T3 hardware root of
trust) stating explicitly what each tier does and does not defend
against, and a canonical, device-signed telemetry bundle that binds the
certified design (G-code SHA-256), the commanded process envelope, the
observed telemetry, and the job outcome into one Ed25519-signed unit
whose hash is committed on-chain.  We implement the T1 tier end-to-end on
real hardware---a Creality Ender-3, OctoPrint on a Raspberry Pi, and
Hyperledger Fabric v3.1.0---and integrate it with an on-chain Bayesian
reputation system, so measured physical adherence becomes a per-actor,
per-dimension trust signal within one atomic transaction.  We evaluate on
\textbf{five real print jobs} spanning three geometries (a complex part, a
calibration cube, and a functional mounting bracket) and three outcome
types (completed, cancelled, and an injected thermal anomaly).  In a
controlled comparison, a $+15\,^{\circ}$C setpoint injection lowered the
derived quality score by 9.3\% while design-hash compliance stayed at
1.0---the bridge separates faithful execution from design substitution
using telemetry alone.  Each job produced a 7--14-kB signed bundle anchored
by a constant 369-B on-chain footprint and updated the operator's Beta
posterior, whose 95\% confidence interval tightened from width 0.684 to
0.561 as attested evidence accumulated; the full attestation pipeline runs
in 1.55\,ms, under 3\% of the ledger-commit cost.  A single flipped sample
was detected by two independent mechanisms (ledger-hash mismatch and
signature failure).  We state the trust boundary honestly: edge signing
makes evidence tamper-evident, but cannot prevent fabrication at source by
a colluding device-key holder---the gap the T2/T3 tiers target.
